\section{Introduction}
Many physical phenomena can be modeled by \glspl{diffEq} together with \glspl{boundaryCondition} and/or
\glspl{initialCondition}. A problem is well posed when a solution exists, is unique, and depends continuously on the prescribed
data. Several related \glspl{diffEq} form a \gls{sysOfDiffEqs}. A \gls{diffEq} involves derivatives of an unknown function.

\Glspl{spatialCoord}\footnote{Also called \glspl{spatialVar}.} and/or time are the most common independent
variables of \glspl{diffEq}. Differential equations may be classified by the number of independent variables as
\gls{ODE} and \gls{PDE}. An \gls{ODE} has just one independent variable, while a \gls{PDE} has two or more independent variables.
In a \gls{linearDiffEq}, the unknown function and its derivatives appear linearly; otherwise, the equation is nonlinear.

\Glspl{diffEq} are defined in a \gls{region}.\footnote{Also called \gls{domain}.} We call end points or surfaces of this region
\glspl{boundary}. A \glsfirst{BVP} consists of a single \gls{diffEq} or a \gls{sysOfDiffEqs} where the solution itself and/or its
derivatives are known on the boundaries. An \gls{IVP}, on the other hand, is a single \gls{diffEq} or a \gls{sysOfDiffEqs}
where the solution itself and/or its derivatives are specified at an initial time, $t = t_0$.

There are two different approaches to solving \glspl{diffEq}, analytical and numerical.\footnote{They result in \glspl{analSol}
or \glspl{exactSol} and \glspl{numSol} or \glspl{approxSol}, respectively.} A real world problem is usually too complicated to be
solved analytically, so we have to resort to numerical approaches. In this report we introduce one of the most famous of these,
\gls{FEM}. \\

In \cref{cha:Problem_Desc} we define a well-known one-dimensional bar problem. Then in \cref{cha:Problem_Formulation} we derive
the governing equations for this problem, before normalizing this in \cref{cha:Problem_Normalization}. The bar problem above is
not fully specified. In \cref{cha:Analytical_Sol} we discuss some problem variations for which the \glspl{analSol} are available.
Finally, \cref{cha:Analytical_Summary} gathers the analytical solutions for all discussed variants.
 
\Cref{cha:FEM_Approach} introduces \gls{FEM} and applies it to the defined problem. In particular, we examine two different choices
of \glspl{basisFunction}.

We are naturally interested in the method's error estimate and the eventual convergence. This is studied in
\Cref{cha:Error_Estimates} for a slightly simplified version of the bar-problem. After deduction of the element matrix with
different approaches in \cref{cha:Results}, we then present results of model application to some problem variations. Finally,
by comparing the FEM with analytical solutions in each case, we investigate if the error estimate in \cref{cha:Error_Estimates}
holds. We conclude this report by a short summary.	\\

A complete list of the terms and descriptions used in this document is provided in Appendix \cref{cha:Glossary_of_Terms}.
Appendix \ref{cha:MatLab_Code} contains the MATLAB code implementing both solution methods, visualization, and error estimation.
%
